\chapter{Состояние методов индексирования многомерных данных}
\label{ch:sota}

\section{Задача доступа к многомерным данным}

Пусть данные организованы в $D$ измерений $H_1,\dots,H_D$, каждое из которых
представляет собой иерархию элементов. Записи адресуются набором
$K=\{h_1,\dots,h_D\}$, $h_i\in H_i$. Аналитический запрос $Q\{h\}$ отбирает все
записи, ключи которых подчинены заданным элементам иерархий, и агрегирует их.

В \cite{borodin2011thesis} показано, что при традиционном для СУБД
индексировании по составному ключу число просматриваемых диапазонов растёт
экспоненциально с $D$, и предложено отображение иерархии на числовую ось:
каждому элементу $h_i$ ставится в соответствие отрезок $[a_i,b_i]$ так, что
отрезок потомка вложен в отрезок предка, а отрезки несвязанных элементов не
пересекаются. При таком отображении аналитический запрос переходит в
пространственный запрос типа <<все точки внутри $D$-мерного параллелепипеда>>, а
задача сводится к пространственному индексированию.

Дальнейшее изложение опирается на это сведение, но не ограничивается им:
рассматриваемые структуры индексируют произвольные многомерные объекты, для
которых определены операции объединения ключей и проверки пересечения.

\section{Деревья пространственного индексирования}

R-дерево \cite{guttman1984r} представляет собой сбалансированное дерево, узел
которого содержит набор ограничивающих прямоугольников (ключей) вместе со
ссылками на поддеревья. Ключ узла покрывает все ключи его потомков. Поиск
спускается по всем поддеревьям, ключи которых пересекаются с запросом; в отличие
от B-дерева, путей спуска может быть несколько, и стоимость поиска определяется
не высотой дерева, а числом узлов, ключи которых пересекаются с запросом.

Отсюда два критерия качества дерева:

\begin{itemize}
\item \emph{перекрытие} ключей узлов одного уровня~--- чем больше суммарная
      площадь пересечений, тем больше поддеревьев приходится просматривать при
      точечном запросе;
\item \emph{покрытая площадь}~--- чем больше <<пустого>> пространства охватывает
      ключ, тем чаще запрос спускается в поддерево, не содержащее ответа.
\end{itemize}

R*-дерево \cite{beckmann1990rstar} улучшает оба показателя за счёт учёта
периметра при разбиении узла и принудительного перевставления части элементов.
RR*-дерево \cite{beckmann2009rrstar} уточняет критерий выбора поддерева при
вставке. Общее для всех вариантов: качество дерева определяется двумя решениями~---
выбором поддерева при вставке (функция штрафа) и разбиением переполненного узла.

\section{Обобщённое дерево поиска}

Обобщённое дерево поиска \cite{hellerstein1995generalized} абстрагирует
описанную схему. Ядро дерева реализует страничную организацию, спуск, вставку,
расщепление, конкурентный доступ и журналирование; предметная часть выносится в
\emph{класс операторов}, который для \pg{} состоит из функций:

\begin{itemize}
\item \code{consistent}~--- может ли поддерево с данным ключом содержать ответ
      на запрос;
\item \code{union}~--- ключ, покрывающий заданный набор ключей;
\item \code{penalty}~--- штраф за расширение ключа при вставке;
\item \code{picksplit}~--- разбиение переполненного узла на две части;
\item \code{same}~--- совпадение ключей;
\item \code{compress}, \code{decompress}~--- преобразование значения во
      внутреннее представление и обратно;
\item \code{fetch}~--- восстановление исходного значения из ключа, необходимое
      для индексного сканирования без обращения к таблице;
\item \code{distance}~--- расстояние до запроса для поиска ближайших соседей.
\end{itemize}

Функции \code{consistent}, \code{union}, \code{penalty} и \code{picksplit}
задают структуру дерева полностью: подставив в них прямоугольники, получаем
R-дерево; подставив сигнатуры~--- дерево полнотекстового поиска; подставив
диапазоны~--- дерево диапазонов.

В реализации \pg{} узел дерева отождествляется со страницей буферного кэша, что
позволяет работать с данными, превышающими объём оперативной памяти. Алгоритмы
спуска и расщепления реализованы без рекурсии по стеку. Конкурентный доступ
организован по схеме Корнакера \cite{kornacker1997concurrency}: страницы
связаны правыми ссылками, а расщепление узла оформлено так, что поиск,
пришедший на расщеплённую страницу, может дойти до нужных данных по правой
ссылке, не удерживая блокировку родителя. Соответствие ключа родителя ключам
потомков восстанавливается отложенно.

\section{Индексирование в основной памяти и битовые карты}

Для аналитических нагрузок с малым числом различных значений в измерении
эффективны битовые карты; для резидентных в памяти наборов~--- структуры,
оптимизированные под отсутствие страничного обмена. Сравнительный анализ этих
подходов и деревьев выполнен в \cite{borodin2011parallel} и
\cite{borodin2012monograph}. Общий вывод: битовые карты выигрывают при низкой
мощности измерений и статических данных, деревья~--- при высокой мощности и
обновляемых данных. Настоящая работа рассматривает второй случай.

\section{Размерность данных и её следствия}

С ростом $D$ доля объёма $D$-мерного шара в описанном кубе стремится к нулю, а
расстояния между случайными точками концентрируются около среднего. Практическое
следствие для древовидных индексов: ключи узлов начинают перекрываться почти
полностью, и поиск вырождается в просмотр всего дерева. Границы применимости
пространственного индексирования по размерности исследованы в
\cite{borodin2016highd}; там же рассмотрены подходы, снижающие эффективную
размерность.

\section{Верификация индексных структур}

Ошибка в реализации метода доступа проявляется не отказом, а молчаливой потерей
данных: запрос перестаёт находить существующие записи. В
\cite{borodin2016debug} описан подход к отладке индексных структур, основанный
на явной проверке инвариантов. Однако средств такой проверки в составе
промышленных СУБД для обобщённого дерева поиска и обратного индекса до недавнего
времени не существовало.

\section{Нерешённые задачи и постановка задач исследования}

Обзор показывает, что литература сосредоточена на двух решениях, определяющих
качество дерева,~--- выборе поддерева и разбиении узла. За её пределами
остаются:

\begin{enumerate}
\item \emph{Ветвистость узла.} Она принимается за характеристику структуры,
      хотя определяется представлением ключа и наличием в узле служебных
      данных, и потому поддаётся оптимизации (глава~\ref{ch:fanout}).
\item \emph{Построение индекса.} Рассматривается как многократная вставка, хотя
      при построении весь набор данных известен заранее
      (глава~\ref{ch:sorted-build}).
\item \emph{Сопровождение.} Сборка мусора, возврат страниц в оборот и
      конкурентность обсуждаются для B-дерева, но не для обобщённого дерева
      поиска (глава~\ref{ch:maintenance}).
\item \emph{Верификация.} Инварианты структуры не формализованы и не
      проверяются в эксплуатации (глава~\ref{ch:verification}).
\end{enumerate}

Этим определяется постановка задач, сформулированных во введении. Прежде чем
переходить к алгоритмам, в главе~\ref{ch:models} строятся модели, показывающие,
на какие именно параметры следует воздействовать.
